\begin{figure}
\centering
\includegraphics[height=0.9\linewidth,angle=90]{../binder/binder-2026-06-22-eco-gwas-nobb.ipynb/binder/teeplots/viz=subplots+ext=.pdf}

\vspace{-1ex}

% https://github.com/mmore500/oee4/blob/5122731aaa9a9c414472cfff1cce0d9e64dd136c/binder/2026-06-22-eco-gwas-nobb.ipynb
\caption{%
\textbf{In co-evolved ``background'' strain from case study treatment, most sites non-adaptive when tested in self context are also non-adaptive when tested in no-background context.}
\footnotesize
As an additional control to test context-specificity of genome sites in case study co-evolved background strain that detected as adaptive only in competition experiments incorporating the case study strain as biotic background (Figure \ref{fig:ecogwas:sites}), a further single-site knockout screen for adaptive sites was conducted under no-background conditions (to supplement primary control under self-context conditions, see Figure \labelcref{fig:fit:nobbg,fig:fit:bbg}).
Dashed horizontal lines indicate neutral expectation for each treatment (0.25 for co-evolved background treatment, 0.5 for no background treatment).
Annotations mark the 15 sites identified as adaptive in co-evolved context competitions, but not adaptive in self context competitions;
most are also non-adaptive under no-background conditions --- with the exception of sites 1242 and 1245, which incur a small fitness penalty when knocked out in no-background competitions.
Excluding these two sites leaves 40\% (13/32) of detected genotype complexity in the background case study strain as specific to its co-evolved ecological partner (i.e., ``focal'' case study strain).
}
\label{fig:ecogwas-nobb-supp}
\end{figure}
